\subsection{A branched projective example}

In this subsection $X=\mathbb P^n$, with integer $n\geq2$ and homogeneous coordinates $[Z_0:\cdots:Z_n]$. Let $\omega$ be the Fubini--Study form normalized by $\int_X\omega^n=1$, and put $\mathrm{d}V=\omega^n$. The current set $\mathscr G_\omega$, determinant density $k_T$, functionals $D_p$, $\rho_\omega$ and threshold $b_\omega$ have the meanings of \cref{sec:integrable-gauduchon-currents}.

A smooth semipositive maximizer need not be strictly positive, even when the threshold solution is smooth. The following example also shows why smoothness of the density root was asserted only off the degeneracy set.

\begin{example}\label{ex:semipositive-projective-maximizer}
For every $n\geq2$, there is a smooth nonnegative density $f$ on $\mathbb P^n$, vanishing precisely on the coordinate hyperplanes, such that $b_\omega(f)=1$ and $D_{f^{1/n}}$ has a smooth semipositive maximizer strictly positive off those hyperplanes. The threshold solution is smooth, but $f^{1/n}$ is not differentiable at a general point of any coordinate hyperplane. There is no smooth strictly positive maximizing form. The normalized greatest exponential solutions converge uniformly to the normalized threshold solution.
\end{example}
\begin{proof}
Define the holomorphic map $F:\mathbb P^n\to\mathbb P^n$ and the semipositive form $A$ by
\[
 F[Z_0:\cdots:Z_n]=[Z_0^2:\cdots:Z_n^2],\qquad
 A=\tfrac12F^*\omega.
\]
The homogeneous formula
\[
 u=\tfrac12\log\!\left(\sum_{i=0}^n|Z_i|^4\right)
       -\log\!\left(\sum_{i=0}^n|Z_i|^2\right)
\]
defines a smooth function on projective space, and $A=\omega+\ddc u$. This form is semipositive and is strictly positive exactly where every homogeneous coordinate is nonzero. Define
\[
 f=A^n/\omega^n,\qquad T=A^{n-1}.
\]
Closed Stokes gives $\int f\omega^n=\int\omega\wedge T=1$. Thus $T\in\mathscr G_\omega$. Every genuine $\omega$-plurisubharmonic potential has extended mass at most one: each bounded truncation has mass one, and restriction away from its truncation level can only decrease that mass. Consequently an inequality $\mathrm{MA}^{\mathrm{ext}}_\omega(v)\geq c f\omega^n$ forces $c\leq1$. The smooth solution $u$ gives the reverse bound, so $b_\omega(f)=1$.

The determinant formula for smooth semipositive forms gives
\[
 k_T=f^{(n-1)/n},\qquad
 D_{f^{1/n}}(T)=\int_{\mathbb P^n}f\omega^n=1.
\]
By \eqref{eq:integrable-current-supremum} and $\rho_\omega(f^{1/n})=b_\omega(f)^{-1/n}=1$, this is a maximizer. Its strict positivity away from the coordinate hyperplanes follows from that of $A$; on a hyperplane $A$ has deficient rank, and so does its cofactor matrix.

For the density root, use the chart $Z_0=1$ with coordinates $z_1,\ldots,z_n$. Pulling back the Fubini--Study volume by $(z_1,\ldots,z_n)\mapsto(z_1^2,\ldots,z_n^2)$ and including the factor $2^{-n}$ gives
\[
 f(z)=2^n\left(\prod_{i=1}^n|z_i|^2\right)
       \frac{(1+\sum_i|z_i|^2)^{n+1}}
            {(1+\sum_i|z_i|^4)^{n+1}}.
\]
Near a point of $\{z_1=0\}$ with all the other coordinates nonzero, its $n$th root is $|z_1|^{2/n}$ times a smooth strictly positive function. Restriction to a real line transverse to that hyperplane proves nondifferentiability, including the absolute-value case $n=2$. Coordinate symmetry gives the same conclusion on every component. The absence of a smooth strictly positive maximizer follows from \cref{cor:kahler-smooth-dual-attainment}, since $f$ has zeros. Finally $f$ is bounded, so \cref{cor:semipositive-dual-lp-limit} gives the uniform exponential limit.
\end{proof}
